\documentclass[10pt,aspectratio=43]{beamer}

\usepackage[english]{babel}
\usepackage{fontspec}
\setmainfont{Latin Modern Roman}
\setsansfont{Latin Modern Sans}
\setmonofont{Cascadia Code}[Scale=0.88]
\usepackage{amsmath, amssymb}
\usepackage{booktabs}
\usepackage{graphicx}
\usepackage{xcolor}
\usepackage{hyperref}
\usepackage{fontawesome5}

\usetheme{default}
\usecolortheme{default}
\setbeamertemplate{navigation symbols}{}
\setbeamertemplate{caption}[numbered]
\setbeamersize{text margin left=8mm,text margin right=8mm}

\colorlet{projectviolet}{violet!65!black}
\colorlet{projectorange}{orange!75!black}
\definecolor{projectgray}{RGB}{78, 84, 94}
\colorlet{projectlight}{black!3}

\setbeamercolor{structure}{fg=projectviolet}
\setbeamercolor{frametitle}{fg=projectviolet}
\setbeamercolor{title}{fg=projectviolet}
\setbeamercolor{subtitle}{fg=projectgray}
\setbeamercolor{normal text}{fg=black}
\setbeamercolor{block title}{fg=white,bg=projectviolet}
\setbeamercolor{block body}{fg=black,bg=projectlight}
\setbeamercolor{alerted text}{fg=projectorange}
\setbeamercolor{footline}{fg=projectgray}

\setbeamertemplate{frametitle}{%
    \vspace{0.35em}
    \insertframetitle\par
    \vspace{0.20em}
    {\color{projectorange}\rule{0.18\linewidth}{1.2pt}%
    \color{projectviolet!35}\rule{0.82\linewidth}{0.5pt}}
}
\setbeamertemplate{footline}{%
    \hfill
    \usebeamercolor[fg]{footline}
    \scriptsize\insertframenumber/\inserttotalframenumber
    \hspace{0.8em}\vspace{0.45em}
}

\setlength{\parskip}{0.25em}
\setlength{\abovedisplayskip}{0.45em}
\setlength{\belowdisplayskip}{0.45em}

\newcommand{\projecttitle}{Neural Network-Based European Option Pricing}
\newcommand{\projectsubtitle}{under Black-Scholes SDE dynamics}
\newcommand{\projectrepo}{\href{https://github.com/matteogiorgi/nn-option-pricing}{\faGithub\hspace{0.35em}\texttt{nn-option-pricing}}}
\newcommand{\placeholderfigure}[1]{%
    \begin{center}
        \fbox{%
            \begin{minipage}[c][0.36\textheight][c]{0.78\linewidth}
                \centering
                \small Figure not available: \texttt{\detokenize{#1}}
            \end{minipage}
        }
    \end{center}
}
\newcommand{\maybeincludegraphics}[2][]{%
    \IfFileExists{#2}{%
        \includegraphics[#1]{#2}%
    }{%
        \placeholderfigure{#2}%
    }%
}
\newcommand{\takeaway}[1]{%
    \begin{block}{Takeaway}
        #1
    \end{block}
}

\title{\projecttitle}
\subtitle{\projectsubtitle}
\author{Matteo Giorgi \and Lorenzo Ghiotti \and Anne Frances Nalumansi \and Alice Pedron}
\institute{Machine Learning for Finance}
\date{May 30, 2026}

\begin{document}

% ------------------------------------------------------------------
% Title
% ------------------------------------------------------------------

\begin{frame}[plain]
    \vspace*{0.65cm}
    {\fontsize{25}{30}\selectfont\color{projectviolet} Neural Network-Based\par}
    \vspace{-0.03cm}
    {\fontsize{25}{30}\selectfont\color{projectviolet} European Option Pricing\par}
    \vspace{0.32cm}
    {\Large under Black-Scholes SDE dynamics\par}

    \vspace{0.75cm}
    {\large\itshape A supervised learning study of neural networks\par}
    {\large\itshape as pricing surrogate models\par}
    \vspace{0.35cm}
    {\footnotesize\color{projectorange}\projectrepo\par}

    \vfill
    \hfill
    \begin{minipage}{0.42\linewidth}
        \raggedleft
        \small
        Matteo Giorgi\par
        Lorenzo Ghiotti\par
        Anne Frances Nalumansi\par
        Alice Pedron\par
    \end{minipage}
\end{frame}

% ------------------------------------------------------------------
% Speaker 1: motivation and theory
% ------------------------------------------------------------------

\begin{frame}{Research Question}
    \begin{block}{Main question}
        Can a feed-forward neural network accurately approximate the Black-Scholes pricing function for European call options?
    \end{block}

    The project is formulated as a controlled supervised-learning problem:
    \[
        (S_0, K, T, r, \sigma)
        \longmapsto
        C_{BS}(S_0,K,T,r,\sigma).
    \]

    \takeaway{Since the target is known exactly, the approximation error can be measured directly.}
\end{frame}

\begin{frame}{Why This Problem Is Meaningful}
    Option pricing often requires evaluating many contracts across many parameter configurations.

    \begin{itemize}
        \item Closed-form formulas are fast but available only in special cases.
        \item Monte Carlo is flexible but can be expensive when repeated many times.
        \item A trained neural network can produce prices through a single forward pass.
    \end{itemize}

    \vspace{0.2em}
    \takeaway{The neural network is studied as a pricing surrogate: a fast approximation of a known pricing map.}
\end{frame}

\begin{frame}{Black-Scholes Dynamics}
    Under the risk-neutral Black-Scholes model:
    \[
        dS_t = rS_t\,dt + \sigma S_t\,dW_t.
    \]

    The exact terminal distribution is:
    \[
        S_T =
        S_0\exp\left[
            \left(r-\frac{1}{2}\sigma^2\right)T
            +\sigma\sqrt{T}Z
            \right],
        \qquad Z\sim\mathcal{N}(0,1).
    \]

    \takeaway{This exact expression is also useful computationally, because Monte Carlo can sample \(S_T\) directly.}
\end{frame}

\begin{frame}{European Call Option}
    A European call option gives the right to buy the underlying asset at strike \(K\) at maturity \(T\).

    \[
        \text{payoff} = (S_T-K)^+ = \max(S_T-K,0).
    \]

    Under risk-neutral pricing:
    \[
        C_0 = e^{-rT}\mathbb{E}^{\mathbb{Q}}\left[(S_T-K)^+\right].
    \]

    \takeaway{All methods in the project are evaluated by how well they recover this price in the Black-Scholes setting.}
\end{frame}

\begin{frame}{Black-Scholes Formula}
    The analytical call price is:
    \[
        C_{BS}
        =
        S_0N(d_1)
        -
        Ke^{-rT}N(d_2),
    \]
    \[
        d_1 =
        \frac{\log(S_0/K)+(r+\frac{1}{2}\sigma^2)T}{\sigma\sqrt{T}},
        \qquad
        d_2=d_1-\sigma\sqrt{T}.
    \]

    \takeaway{The formula provides the clean target variable used for supervised learning.}
\end{frame}

\begin{frame}{Monte Carlo Benchmark}
    Monte Carlo approximates the risk-neutral expectation by simulation:
    \[
        \widehat{C}_{MC}
        =
        e^{-rT}\frac{1}{M}
        \sum_{i=1}^{M}
        (S_T^{(i)}-K)^+.
    \]

    \begin{itemize}
        \item It is flexible and widely used.
        \item It is stochastic, so estimates fluctuate.
        \item The standard error decreases at rate \(M^{-1/2}\).
    \end{itemize}
\end{frame}

% ------------------------------------------------------------------
% Speaker 2: data and ML methodology
% ------------------------------------------------------------------

\begin{frame}{Synthetic Dataset}
    The dataset is generated by sampling option parameters and computing the analytical Black-Scholes price.

    \begin{center}
        \small
        \begin{tabular}{lll}
            \toprule
            Variable   & Meaning           & Range         \\
            \midrule
            \(S_0\)    & underlying price  & \([50,150]\)  \\
            \(K\)      & strike price      & \([50,150]\)  \\
            \(T\)      & maturity in years & \([0.1,2]\)   \\
            \(r\)      & risk-free rate    & \([0,0.05]\)  \\
            \(\sigma\) & volatility        & \([0.1,0.6]\) \\
            \bottomrule
        \end{tabular}
    \end{center}
\end{frame}

\begin{frame}{Why Synthetic Data?}
    The project does not aim to model observed exchange-traded option quotes.

    \begin{itemize}
        \item Synthetic labels give a known ground truth: \(C_{BS}\).
        \item The experiment is controlled and reproducible.
        \item Real option prices include bid-ask spreads, liquidity, dividends, and implied-volatility effects.
    \end{itemize}

    \begin{block}{Methodological choice}
        Real market data would change the research question from approximating Black-Scholes to modeling market prices.
    \end{block}
\end{frame}

\begin{frame}{Feature Engineering: Moneyness}
    The final feature set includes moneyness:
    \[
        \frac{S_0}{K}.
    \]

    The neural network receives:
    \[
        (S_0,K,T,r,\sigma,S_0/K).
    \]

    \begin{itemize}
        \item \(S_0/K<1\): out-of-the-money region.
        \item \(S_0/K\approx 1\): at-the-money region.
        \item \(S_0/K>1\): in-the-money region.
    \end{itemize}
\end{frame}

\begin{frame}{Supervised Learning Setup}
    Each observation has the form:
    \[
        x_i=(S_0,K,T,r,\sigma,S_0/K)_i,
        \qquad
        y_i=C_{BS,i}.
    \]

    The model learns a function:
    \[
        f_\theta(x_i)\approx y_i.
    \]

    Inputs and targets are standardized during training, while metrics are reported in original price units.
\end{frame}

\begin{frame}{Neural Network Architecture}
    The selected model is a feed-forward multilayer perceptron.

    \begin{center}
        \small
        \begin{tabular}{ll}
            \toprule
            Component      & Final choice               \\
            \midrule
            Input features & \(S_0,K,T,r,\sigma,S_0/K\) \\
            Hidden layers  & fully connected MLP        \\
            Activation     & SiLU                       \\
            Output         & one scalar price           \\
            \bottomrule
        \end{tabular}
    \end{center}

    \takeaway{The architecture is simple but appropriate for smooth tabular regression.}
\end{frame}

\begin{frame}{Training and Metrics}
    The training procedure uses:
    \begin{itemize}
        \item 70\% / 15\% / 15\% train-validation-test split;
        \item MSE loss;
        \item Adam optimizer;
        \item early stopping on validation loss.
    \end{itemize}

    Main metrics:
    \[
        \mathrm{MAE},\qquad
        \mathrm{RMSE},\qquad
        R^2,\qquad
        \mathrm{MAPE}\text{ for prices }>1.
    \]
\end{frame}

% ------------------------------------------------------------------
% Speaker 3: implementation and experiments
% ------------------------------------------------------------------

\begin{frame}{Software Architecture}
    The code is organized as a reusable Python package plus command-line experiment scripts.

    \begin{center}
        \scriptsize
        \begin{tabular}{ll}
            \toprule
            Module                     & Responsibility            \\
            \midrule
            \texttt{black\_scholes.py} & analytical pricing        \\
            \texttt{dataset.py}        & synthetic data generation \\
            \texttt{model.py}          & neural network            \\
            \texttt{training.py}       & training and prediction   \\
            \texttt{monte\_carlo.py}   & simulation benchmark      \\
            \texttt{svr.py}            & classical ML baseline     \\
            \bottomrule
        \end{tabular}
    \end{center}

    \takeaway{Python orchestrates optimized numerical backends; the heavy computations are not pure Python loops.}
\end{frame}

\begin{frame}{Experiment Pipeline}
    The main experiment follows a deterministic workflow:

    \begin{enumerate}
        \item generate synthetic contracts;
        \item compute analytical prices;
        \item split and scale the data;
        \item train the neural network;
        \item evaluate against \(C_{BS}\);
        \item save metrics, plots, configuration, and artifacts.
    \end{enumerate}

    \takeaway{The experiment is reproducible from the command line, not tied to an interactive notebook.}
\end{frame}

\begin{frame}{Final Experiment Configuration}
    The final neural-network run uses:

    \begin{center}
        \small
        \begin{tabular}{ll}
            \toprule
            Setting           & Value                      \\
            \midrule
            Dataset size      & 100,000 contracts          \\
            Feature set       & base variables + moneyness \\
            Activation        & SiLU                       \\
            Maximum epochs    & 200                        \\
            Batch size        & 1024                       \\
            Monte Carlo paths & 50,000                     \\
            \bottomrule
        \end{tabular}
    \end{center}
\end{frame}

\begin{frame}{Model Selection Experiments}
    Intermediate experiments justify the final configuration.

    \begin{center}
        \scriptsize
        \begin{tabular}{lrrr}
            \toprule
            Experiment       & MAE     & RMSE    & MAPE     \\
            \midrule
            Base + ReLU      & 0.17269 & 0.22973 & 1.8687\% \\
            Moneyness + ReLU & 0.15417 & 0.20930 & 1.6670\% \\
            Base + SiLU      & 0.10775 & 0.15817 & 1.2488\% \\
            Moneyness + SiLU & 0.10254 & 0.15169 & 1.3213\% \\
            \bottomrule
        \end{tabular}
    \end{center}

    \takeaway{The combined moneyness + SiLU setup gives the strongest MAE/RMSE candidate in the intermediate experiments.}
\end{frame}

\begin{frame}{Runtime Benchmark}
    Runtime is measured after training, for 1,000 pricing queries.

    \begin{center}
        \scriptsize
        \begin{tabular}{lrr}
            \toprule
            Method                    & Time      & Options/s \\
            \midrule
            Black-Scholes formula     & 0.00091 s & 1,101,095 \\
            Neural inference          & 0.03190 s & 31,351    \\
            Monte Carlo, 20,000 paths & 1.28235 s & 780       \\
            \bottomrule
        \end{tabular}
    \end{center}

    \takeaway{The analytical formula is fastest, but neural inference is about 40 times faster than Monte Carlo in this benchmark.}
\end{frame}

\begin{frame}{SVR Baseline}
    Support Vector Regression is included as a classical machine-learning reference.

    \begin{center}
        \small
        \begin{tabular}{lrr}
            \toprule
            Metric  & Mean     & Std      \\
            \midrule
            MAE     & 0.15088  & 0.00732  \\
            RMSE    & 0.23100  & 0.00592  \\
            \(R^2\) & 0.999905 & 0.000005 \\
            MAPE    & 1.8253\% & 0.2422\% \\
            \bottomrule
        \end{tabular}
    \end{center}

    \small The benchmark is intentionally reduced to 5,000 samples over three seeds because RBF kernel methods scale less favorably.
\end{frame}

% ------------------------------------------------------------------
% Speaker 4: results and conclusions
% ------------------------------------------------------------------

\begin{frame}{Final Neural Network Results}
    \begin{center}
        \small
        \begin{tabular}{lr}
            \toprule
            Metric              & Value    \\
            \midrule
            MAE                 & 0.04290  \\
            RMSE                & 0.06208  \\
            \(R^2\)             & 0.999993 \\
            MAPE, prices \(>1\) & 0.4803\% \\
            \bottomrule
        \end{tabular}
    \end{center}

    \takeaway{The final model approximates the Black-Scholes pricing function with very low error on held-out synthetic contracts.}
\end{frame}

\begin{frame}{Predicted vs Analytical Prices}
    \centering
    \maybeincludegraphics[width=0.60\linewidth]{../results/final/figures/true_vs_predicted.png}

    \vspace{0.1em}
    \small Predictions are concentrated near the identity line, consistently with \(R^2 \approx 1\).
\end{frame}

\begin{frame}{Error Diagnostics}
    \centering
    \maybeincludegraphics[width=0.66\linewidth]{../results/final/figures/error_distribution.png}

    \vspace{0.2em}
    \small The error distribution checks whether low aggregate metrics hide large systematic deviations.
\end{frame}

\begin{frame}{Noisy Target Robustness}
    Training labels are perturbed, while evaluation remains against clean Black-Scholes prices.

    \begin{center}
        \scriptsize
        \begin{tabular}{lrr}
            \toprule
            Noise level & NN MAE  & SVR MAE \\
            \midrule
            0\%         & 0.10254 & 0.15088 \\
            1\%         & 0.11084 & 0.19284 \\
            5\%         & 0.17827 & 0.56185 \\
            \bottomrule
        \end{tabular}
    \end{center}

    \takeaway{Increasing label noise worsens performance, but the neural network remains reasonably stable under mild perturbations.}
\end{frame}

\begin{frame}{Interpretation and Limits}
    The results support a clear but conditional conclusion.

    \begin{itemize}
        \item Strong accuracy holds inside the sampled Black-Scholes domain.
        \item The model is a surrogate, not a replacement for the analytical formula.
        \item Synthetic data are appropriate for this research question.
        \item Real market data would define a different problem.
        \item No-arbitrage constraints are not explicitly enforced.
    \end{itemize}
\end{frame}

\begin{frame}{Final Takeaway}
    \begin{block}{Answer}
        A feed-forward neural network can accurately approximate European call prices generated by the Black-Scholes formula in a controlled synthetic setting.
    \end{block}

    \vspace{0.3em}
    Future work:
    \begin{itemize}
        \item Greeks by automatic differentiation;
        \item stochastic-volatility models;
        \item path-dependent or American options;
        \item real market data as a separate research question.
    \end{itemize}

    \vfill
    \centering
    {\footnotesize\color{projectorange}\projectrepo}
\end{frame}

\end{document}
